\documentclass[10pt,a4paper]{amsart}
\usepackage[margin=19mm]{geometry}
\usepackage{amsmath,amssymb,mathtools}
\usepackage[scr=rsfs,cal=euler]{mathalfa}
\usepackage{xfrac}
\usepackage[unicode,hidelinks,bookmarks=false]{hyperref}
\newcommand{\ie}{i.e.\ }
\newcommand{\cf}{cf.\ }
\newcommand{\resp}{respectively\ }
\newcommand{\Subsection}{\subsection}
\providecommand{\nobreakdash}{\nobreak\hbox{-}}
\setlength{\emergencystretch}{2em}
\multlinegap0pt
\usepackage{bm}
\usepackage{bbm}
\usepackage{dsfont}
\usepackage{xspace}
\usepackage{cancel}
\usepackage{mathtools}
\usepackage{amsthm}
\makeatletter
\@ifundefined{theorem}{\newtheorem{theorem}{Theorem}}{}
\makeatother
\newcommand{\RSpec}{\operatorname{RSpec}}
\newcommand{\Spec}{\operatorname{Spec}}
\newcommand{\nil}{\operatorname{nil}}
\title[Coherent Conditions: Algebraic Geometry for Arbitrary Classe]{Coherent Conditions: Algebraic Geometry for Arbitrary Classes of Algebras}
\author{K. R. van Nispen}
\date{Source: arXiv:2510.23647v1 (CC BY 4.0)}
\begin{document}
\maketitle

\noindent\emph{Source and licence.} Coherent Conditions: Algebraic Geometry for Arbitrary Classes of Algebras. Version arXiv:2510.23647v1 (CC BY 4.0), distributed under the Creative Commons Attribution 4.0 International licence, as recorded in the arXiv licence metadata for this exact version and shown on the abstract page https://arxiv.org/abs/2510.23647v1. Full rights notice: licenses/arxiv-2510.23647-CC-BY-4.0.txt.\par\smallskip

\noindent\textit{Editorial scope.} Counted unit: the Theorem whose source label is \texttt{theo:radfundclassisprevar} in the article (source0.tex lines 470--474, source label \texttt{theo:radfundclassisprevar}), delivered together with the article's own complete original proof of it (source0.tex lines 475--497, one \texttt{proof} environment of 23 lines). No mathematics is written, repaired or re-derived: statement and proof are the article's own text line for line. Required context retained uncounted: theo:radiscc (lines 446--448). Every definition, display and label of those lines is retained unchanged. Omitted from this package and not redistributed: 1--445, 449--469, 498--1179. Nothing inside a retained excerpt is omitted or altered: every retained body is delivered exactly as the article's lines are written, so no omission receipt of any kind accompanies this package. Citations: the retained excerpts carry no citation command at all, so no bibliography entry of the article is transcribed and no reference is redistributed. Reference adaptation: none was needed and none was made; every reference inside the delivered unit resolves inside it, so no unresolved reference or citation is produced. Printed slips kept literal: no printed word, symbol, hypothesis or number of the counted statement or of its proof was altered. Layout and class: the article's own class is \texttt{article} with packages including graphicx, hyperref, amsmath, amssymb, amsthm, tikz, biblatex, geometry, mathpazo, newtxmath; the wrapper is the lane's pinned \texttt{amsart} editorial wrapper at 10pt on A4, which restates the theorem-like environments the retained text needs and adds the article's own macro definitions that the retained text needs (\texttt{\textbackslash RSpec}, \texttt{\textbackslash Spec}, \texttt{\textbackslash nil}), with extra wrapper packages (bm, bbm, dsfont, xspace, cancel, mathtools). The delivered numbering is the unit's own. Assets: no retained span contains a figure environment, a graphics-inclusion command, a PDF-page inclusion, a table or a TikZ picture; no image file is copied into this package, the package declares no asset dependency and no image file is redistributed with it. Licence: version arXiv:2510.23647v1 of this article is distributed under the Creative Commons Attribution 4.0 International licence, as recorded in the arXiv licence metadata for this exact version and shown on the abstract page https://arxiv.org/abs/2510.23647v1. A full rights notice is delivered at licenses/arxiv-2510.23647-CC-BY-4.0.txt. Characters: every retained excerpt is pure ASCII, so the delivered engine, XeLaTeX with its default Latin Modern text font, reports no missing character.
\begin{theorem} \label{theo:radiscc}
    For any coherent condition $\varphi$ on $V$ there is a coherent condition $\sqrt{\varphi}$ on $V$ such that $\Spec^{\sqrt{\varphi}}\mathbf{A} = \RSpec^\varphi\mathbf{A}$.
\end{theorem}

\begin{theorem} \label{theo:radfundclassisprevar}
    Let $\varphi = \Phi_V(K)$ be a coherent condition.
    Then $F(\sqrt{\varphi}) = ISP(K)$.
    In other words, $\sqrt{\varphi} = \Phi_V(ISP(K))$, and so $\sqrt{\sqrt{\varphi}} = \sqrt{\varphi}$.
\end{theorem}

\begin{proof}
    We will show this by using a bi-inclusion.
    For the first direction, let $\mathbf{A} \in F(\sqrt{\varphi})$.
    Consider the natural map 
    $$\alpha \colon \mathbf{A} \rightarrow \prod_{\theta \in \Spec^\varphi\mathbf{A}}\mathbf{A}/\theta$$
    given by 
    $$a \mapsto (a/\theta)_{\theta \in \Spec^\varphi\mathbf{A}}$$
    This has kernel
    $$\bigcap\Spec^\varphi\mathbf{A} = \nil^\varphi\mathbf{A} = \Delta_A$$
    hence $\alpha$ is an embedding.
    Further, as $\mathbf{A}/\theta$ is in $IS(K)$ for all $\varphi$-distinguished $\theta$ we have
    $$\mathbf{A} \cong \operatorname{im}\alpha \leq \prod_{\theta \in \Spec^\varphi\mathbf{A}}\mathbf{A}/\theta \in ISP(K)$$
    $$\implies \mathbf{A} \in ISP(K)$$
    Hence $F(\sqrt{\varphi}) \subseteq ISP(K)$.\par
    Conversely, by Theorem \ref{theo:radiscc} it suffices to show that $F(\sqrt{\varphi})$ contains $P(K)$.
    Thus let $\mathbf{A}_i \in K$ be a family of algebras indexed by $I$, and define $\mathbf{A} = \prod_{i \in I}\mathbf{A}_i$.
    Let $\theta_i$ be the kernel of projection map $\pi_i \colon \mathbf{A} \rightarrow \mathbf{A}_i$.
    It should be clear that $\bigcap_{i \in I} \theta_i = \Delta_A$.
    But obviously $\theta_i \in \Spec^\varphi\mathbf{A}$, so
    $$\Delta_A = \bigcap_{i \in I} \theta_i \in \RSpec^\varphi\mathbf{A}$$
    Hence $P(K) \subseteq F(\sqrt{\varphi})$.\par
    That $\sqrt{\sqrt{\varphi}} = \sqrt{\varphi}$ is then translates to the fact that $ISP(ISP(K)) = ISP(K)$.
\end{proof}

\end{document}
